% !TeX root = ../report.tex
% !Mode:: "TeX:UTF-8"
% 哈尔滨工业大学本科毕业设计（论文）开题报告正文，detect 阶段模板。
% 封面见 ../front/coverart.tex，主文件见 ../report.tex。

\section{课题来源及研究的目的和意义}

\subsection{课题来源}

本课题依托围绕软件技术债修复与测试用例演化开展的科研项目，该项目主要关注软件在功能修改、缺陷修复或依赖调整过程中，测试用例如何随生产代码的变化而演化，并探索利用代码变更信息、测试意图和大语言模型（Large Language Model，LLM）辅助测试维护的方法。在此基础上，本课题进一步将研究场景扩展至一般拉取请求（Pull Request，PR），重点关注测试用例更新前的目标识别问题。


\subsection{研究目的和意义}

随着软件功能修改、缺陷修复和依赖调整等变更，已有测试用例也需要随生产代码的变化进行相应更新。现有 Coding Agent 通常优先完成生产代码修改，对已有测试用例的同步维护关注不足。若仅依据已经完成的代码差异更新测试，生成的测试容易依赖具体实现细节，削弱测试对预期行为的独立验证能力。另一方面，现有大语言模型研究更多关注新测试生成；若持续新增测试文件，可能造成测试冗余和规模膨胀，增加测试套件的理解与维护成本。

而且，若测试用例未随需求同步更新，Coding Agent 可能依据过时的测试验证新实现，导致验证结论失真；若测试更新完全依赖已经完成的源代码修改，又容易复现实现中的假设与错误，形成循环验证，降低测试发现缺陷的能力。所以，测试用例更新应以开发任务所描述的预期行为为主要依据，在目标定位阶段与具体实现结果保持信息隔离，从而增强测试验证的独立性和可信度。

随着 Coding Agent 在软件开发中的应用不断扩大，代码生成与修改效率持续提高，测试维护与验证能力不足的问题也日益突出。因此，有必要及时研究面向 Coding Agent 的测试用例更新方法，建立可靠的目标定位与评价机制，避免测试维护滞后于生产代码自动化，进而制约智能化软件开发的质量与效率。

因此，相较于无差别生成新测试，更合理的方式是根据开发任务识别并更新已有测试。在这一过程中，首先需要解决的问题是：给定一个开发任务，哪些已有测试需要随之更新？由于开发任务通常通过 Issue 或 PR 的标题和正文等自然语言进行描述，本课题聚焦于根据自然语言描述和当前代码仓库，定位并选择需要更新的已有测试文件。

本课题关注如何利用开发者在 PR 中提供的自然语言描述，识别代码仓库中需要随软件变更进行更新的已有测试用例。以 Scrapy 项目的实际 PR 为例，Scrapy 是一个基于 Python 的开源网络爬虫与数据抓取框架，主要用于网站数据采集和结构化信息提取。 在 Scrapy 的 PR 7846 中，开发者针对 Twisted 中 IDNA 辅助函数的兼容性问题进行了调整，并在 PR 的标题和正文中描述了此次变更的背景及预期行为。
随着生产代码的修改，该 PR 同时更新了下载器核心测试文件 \texttt{tests/test\_core\_downloader.py} 和共享 HTTP 下载测试基类 \texttt{tests/utils/bases/download\_handlers\_http.py}，以使已有测试适应软件行为的变化。
这一实例表明，PR 标题和正文中的开发者描述能够反映软件变更的背景、意图和预期行为，从而为识别需要随软件变更进行更新的测试用例提供依据。

软件不断演化的过程中，测试用例也需要随功能、接口和依赖关系的变化进行相应维护，以保证其能够持续反映软件的预期行为。
然而，在规模较大的代码仓库中，一次软件变更可能同时影响直接测试、共享测试基类以及通过调用关系间接受影响的测试，仅依靠开发人员人工分析代码和测试之间的关联，不仅需要投入较多时间，也容易遗漏需要同步更新的测试用例。
因此，如何准确识别受软件变更影响的已有测试，是测试用例更新过程中需要首先解决的问题。

PR 中的自然语言描述通常包含开发者对变更背景、修改原因和预期行为的说明，能够提供仅从代码结构和文件关系中难以直接获得的变更意图信息。
利用这些信息辅助识别需要更新的测试用例，有助于缩小测试维护过程中的人工排查范围，为后续测试用例更新提供更加明确的目标。
同时，研究软件变更意图与已有测试用例之间的关联，也有助于更充分地利用开发过程中已有的自然语言信息，为提高测试维护的效率和准确性提供基础。

基于上述分析，本课题将测试用例更新前的目标识别作为研究重点，旨在从已有测试中准确识别需要随软件变更进行维护的测试用例。准确识别需要更新的测试用例，不仅能够缩小开发人员的人工排查范围，减少遗漏关键测试的风险，也能够为后续测试修改提供更加明确的依据。识别范围过大将增加额外的审查成本，而识别范围过小则可能遗漏需要同步维护的测试。因此，本课题重点关注测试用例更新之前的识别环节，探索如何充分利用 PR 中的自然语言变更信息和代码仓库上下文，提高测试用例识别的准确性和完整性。

近年来，Coding Agent 对自然语言需求和代码仓库上下文的理解能力不断增强。它不仅能够根据开发者提供的自然语言描述理解软件变更的目的和预期行为，还能够分析相关模块、调用关系以及代码与测试之间的联系。当功能、接口或依赖发生变化时，已有测试用例可能因测试输入、模拟对象或断言仍对应旧行为而失效或不再适用。借助 Coding Agent，可以将自然语言表达的变更需求与仓库中的已有测试联系起来，定位可能受到代码变更影响、需要同步更新的测试文件，为后续测试维护缩小排查范围。因此，本课题重点关注测试用例更新之前的识别环节，探索如何充分利用 PR 中的自然语言变更信息和代码仓库上下文，提高测试用例识别的准确性和完整性。

然而，现有大语言模型尚不能仅依据开发任务描述和修改前的代码仓库，稳定、准确地定位需要更新的已有测试用例。前期实验表明，即使借助 SDD 工作流组织需求分析，模型仍可能产生无关候选、遗漏实际目标，并且在不同 PR 上的表现存在明显波动。这说明开发任务与测试文件之间仍缺少可靠的关联和选择机制，也表明有必要通过受控实验系统分析模型的定位缺陷，并研究更加准确的测试更新目标定位方法。

围绕上述研究问题和技术思路，本课题主要从实验流程构建、不同方法比较和影响因素分析三个方面展开研究。第一，建立从 PR 样本采集到文件级评价的可复现实验流程，检验自然语言能否支持已有测试文件定位。第二，基于 OpenSpec、Spec Kit 和 cc-sdd 三种规范驱动开发（Specification-Driven Development，SDD）流程，研究自然语言需求驱动的测试用例更新目标定位方法。第三，分析自然语言质量、代码与测试的关联方式以及修改规模对误检和漏检的影响，为后续测试更新方法提供依据。

本课题具有较强的理论价值和实际应用价值。从研究意义看，本课题以 OpenSpec、Spec Kit 和 cc-sdd 等 SDD 工具为基础，探索如何将自然语言描述的软件变更意图逐步转化为测试用例更新方案，并研究测试目标定位在需求理解与测试修改之间的衔接作用。从工程意义看，准确定位需要更新的已有测试文件能够缩小 Coding Agent 的搜索范围，为后续修改测试输入、模拟对象和断言提供明确依据，最终为实现自然语言驱动的测试用例自动更新奠定基础。

\section{国内外在该方向的研究现状及分析}

\subsection{基于规则和程序关系的测试选择方法}

传统的测试选择与变更影响分析主要依据确定性规则和程序关系，识别可能受到代码变更影响的测试。常见线索包括文件命名约定、目录对应关系、生产代码与测试代码的静态依赖、测试覆盖信息以及历史共同修改记录。此类方法可以根据代码变更快速缩小测试范围，结果容易解释，并且不需要大规模语言模型推理。

在测试共同演化研究中，生产代码实体与测试实体之间的历史链接是重要信息。Liu 等\cite{liu2025cote}提出 COTE，将链接分析与预训练语言模型结合，用于预测代码与测试之间的共同演化关系。这说明历史修改关系能够为测试选择提供证据。但规则和程序关系方法通常依赖可获得且准确的代码差异、覆盖数据或历史链接；当检测发生在变更实现之前，或者测试通过公共基类、夹具和间接调用覆盖功能时，单一结构关系难以完整反映开发者的变更意图。

\subsection{基于语义表示与检索的测试定位方法}

为弥补程序关系难以直接表达软件变更意图的问题，语义表示方法尝试将自然语言和代码映射到统一的表示空间，以衡量二者之间的语义关联。Feng 等\cite{feng2020codebert}提出 CodeBERT，通过联合学习编程语言和自然语言表示支持自然语言代码检索等任务。Guo 等\cite{guo2022unixcoder}提出 UniXcoder，利用代码结构和跨模态信息学习统一的代码表示。这些模型为比较 PR 自然语言与测试代码的语义相关性提供了基础。

检索方法可以先使用关键词、BM25 或向量相似度缩小候选集合，再利用模型进行语义判断。IssueExec\cite{liu2026issueexec} 使用测试作为软件工程问题与生产代码之间的中间层：利用历史提交和共同变更实体补充测试表示，以 BM25 检索候选测试，再由大语言模型选择与 issue 相关的测试，最终结合执行轨迹定位生产代码。本课题拟借鉴其“问题文本到相关测试”的检索阶段，并将测试函数结果归并为文件路径。

语义检索能够处理自然语言与测试代码之间的词汇差异，但“语义相关”并不等于“需要更新”。某个测试可能与变更主题相关却无需修改，也可能因为间接依赖而需要更新但文本相似度较低。因此，检索结果必须在本课题的更新文件标签和评价协议下重新验证。

\subsection{基于大语言模型的测试生成与更新方法}

大语言模型能够结合文本意图和代码上下文生成、改进或修复测试。Nan 等\cite{nan2025intut}提出 IntUT，通过测试输入、模拟对象行为和预期结果等测试意图引导大语言模型生成单元测试。UnitCon\cite{jang2025unitcon} 面向 Java 运行时异常合成有针对性的单元测试。Schaefer 等\cite{schaefer2024empirical}对大语言模型自动生成单元测试进行了实证评价，Alshahwan 等\cite{alshahwan2024testgenllm}研究了工业环境中的大语言模型单元测试改进。

对于已有测试的修复与更新，Fatima 等\cite{fatima2024flakyfix}提出 FlakyFix，使用大语言模型预测不稳定测试修复类别并修复测试代码。Tian 等\cite{tian2026mumutestup}提出 MuMuTestUp，利用变异分析、覆盖分析和语义检索智能体生成反馈，迭代更新已有测试。这些研究表明，大语言模型可以理解测试意图并生成修改，但多数工作假设目标类、目标方法或目标测试已经确定，重点在于“怎样生成或修改”，没有完整解决“应先修改哪个已有测试文件”。

\subsection{基于 SDD 的规划方法}

为在代码实现之前更加系统地理解和规划自然语言需求，规范驱动开发（Specification-Driven Development，SDD）通常将需求逐步转化为规格、设计和任务等中间产物，再进入具体实现阶段。根据 OpenSpec 官方项目文档\cite{openspecOfficial}，其流程使用 proposal、specs、design 和 tasks 等产物描述变更；GitHub Spec Kit 官方项目\cite{speckitOfficial}将 specification 逐步展开为 plan 和 tasks；cc-sdd 官方项目\cite{ccsddOfficial}采用 requirements、design 和 tasks 等阶段式文档组织需求与实现关系。虽然三种工具在产物名称和具体流程上存在差异，但其共同特点是利用自然语言需求和代码仓库上下文逐步形成结构化的规划结果，并在规划过程中记录可能受到影响的模块和测试任务。

与直接根据关键词或语义相似度检索测试相比，SDD 的特点在于先理解软件变更需求，再结合仓库结构分析其可能产生的影响。这一过程能够形成显式的规格、设计和任务等规划产物，为分析自然语言变更意图与已有测试之间的联系提供更多上下文信息。但不同工具的产物格式、规划粒度和路径表达不同，需要统一解析后才能比较；较长的规划链也可能扩大候选范围或引入与任务无关的推断。

\subsection{现有方法存在的问题与本课题切入点}

现有工作已经从程序关系、语义检索、LLM 和需求规划等角度探索了代码与测试之间的关联，但对于如何从软件变更描述中识别需要同步更新的已有测试，仍存在若干问题。

\begin{enumerate}
  \item \textbf{大语言模型难胜任，目标定位仍欠缺。}现有大语言模型在 OpenSpec、Spec Kit 和 cc-sdd 等 SDD 工作流的支持下，虽然能够将自然语言需求组织为需求说明、设计方案和任务列表，但这些规划过程并非专门面向已有测试用例更新。预实验结果表明，模型仍会产生较多无关候选或遗漏实际目标，且在不同 PR 上的表现存在较大差异。因此，规划内容的完整性尚不能直接转化为测试更新目标定位的准确性和完整性。

  \item \textbf{输出粒度有差异，实验结果难比较。}已有方法可能输出测试类、测试函数、生产代码位置或文本形式的修改建议，而本课题关注需要更新的已有测试文件。若不统一输出粒度和路径规则，不同方法的结果就无法进行一致的定量评价。

  \item \textbf{更新定位有不足，实际效果难验证。}现有大语言模型研究更多关注如何生成、改进或修复测试，通常假设目标类、目标方法或目标测试已经确定。对于如何根据 PR 自然语言从整个仓库中定位待更新测试文件，仍缺少统一数据和评价协议下的系统比较。
\end{enumerate}

总体来看，本课题针对现有研究以 PR 自然语言描述和 base 版本仓库为输入，以需要更新的已有测试文件为定位目标。在统一信息边界、输出粒度和评价方式的基础上，通过比较不同 SDD 流程及检索方法的定位效果，进而分析大语言模型产生误检和漏检的原因，并进一步研究自然语言需求表示、代码与测试关联分析以及候选文件选择方法。最终目标是提高大语言模型根据自然语言描述定位测试更新目标的准确性和完整性，为后续测试用例自动更新提供可靠依据。

\section{主要研究内容}

\subsection{总体任务与研究内容}

本课题研究的问题是：给定 PR 标题、PR 正文和 base 版本仓库，定位并选择可能需要随该 PR 更新的已有测试文件。设 PR $p$ 的输入为 $I_p=\{T_p,B_p,R_p^{base}\}$，方法最终输出的候选测试文件集合为 $C_p$，历史实际修改的参考测试文件集合为 $G_p$。$G_p$ 只在评价阶段使用，不进入方法输入。

样本采集和质量检查属于实验前期准备工作，其规则和已完成情况统一放在第五章说明。围绕“从自然语言到文件级候选集合”这条主线，后续研究包含两个相互衔接的部分：首先理解 PR 自然语言和 base 仓库，发现可能需要更新的原始测试目标；随后把原始目标转换为有效的文件级候选，并评价候选集合与历史修改之间的差异。前者回答“可能需要更新哪些测试”，后者回答“哪些定位结果能够作为最终候选、候选效果如何”。

\subsection{自然语言驱动的测试更新目标检测}



测试更新目标检测首先需要解决自然语言需求与测试代码之间缺少直接对应关系的问题。PR 标题和正文通常描述功能目的、兼容性要求或用户可见行为，而测试文件可能使用不同的术语，并通过公共基类、夹具、继承关系或间接调用覆盖相关功能。方法既要理解开发者希望改变什么，又要结合 base 仓库判断哪些生产模块和已有测试可能受到影响。

同时还需要解决检测边界和方法可比性问题。历史 PR 已经包含 head 代码与测试修改，如果方法读取这些内容，就不再是在预测更新目标；如果三种 SDD 流程获得不同的自然语言、仓库版本、模型配置或权限，观察到的差异也无法归因于方法本身。因此，需要把历史 PR 构造成参考答案不可见的统一检测任务，并在相同条件下执行 OpenSpec、Spec Kit、cc-sdd 及后续检索对照。



具体研究将 PR 标题、正文和 base 仓库组织为统一输入，独立封存 head 代码、历史测试修改和其他方法产物。三种 SDD 流程只执行规划和 detect，通过需求、规格、设计与任务等中间产物逐步解释变更意图、分析受影响模块，并提出可能需要更新的原始测试目标。实验同时保存规划文档、原始路径或测试符号、判断理由以及成功、失败、中断和显式重试状态，使检测过程能够追溯。

对照方法拟适配 IssueExec 的测试检索阶段，以相同 PR 自然语言和 base 测试集合进行候选检索。通过比较 SDD 规划推理与检索方法的输出，研究显式需求规划是否有助于发现自然语言中未直接出现、但通过代码关系受到影响的测试文件。

\paragraph{理想结果}

理想情况下，方法能够在完全不可见 head 修改答案的前提下，正确理解 PR 描述的变更目的和预期行为，找到所有确实需要更新的已有测试文件，同时避免把仅在主题上相关但无需修改的测试大量纳入候选。对于每个原始目标，方法还应给出可解析的仓库路径及明确理由，说明该测试与自然语言需求、受影响代码或测试依赖之间的联系，并生成完整、可复核且无运行异常的规划产物。

\subsection{文件级目标选择与效果评价}



文件级目标选择首先需要解决异构输出无法直接比较的问题。三种 SDD 工具和检索方法可能输出绝对路径、相对路径、测试类、测试函数或自然语言形式的建议，其中还可能包含重复路径、无法解析的符号以及 base 仓库中不存在的文件。如果直接统计这些原始产物，格式差异会被误认为检测能力差异。因此，需要建立统一的文件级表示和选择规则。

效果评价还需要解决评价口径与结果解释问题。历史实际修改的测试文件只是可观测的代理标签，未修改文件不一定与变更无关；不同 PR 的参考文件数量差异较大，少数多文件 PR 可能主导总体结果；上游中断、空输出和方法正常完成但未命中的情况也不能混为一谈。因此，需要同时评价候选准确程度、参考文件覆盖程度和运行完整性，并通过人工案例解释指标背后的原因。


具体研究从规划产物和检索结果中提取原始测试目标，将路径统一转换为 base 仓库相对路径，检查目标是否存在、是否属于测试文件，并删除重复项。目标定位负责发现可能相关的原始测试目标，目标选择则依据路径可解析性、文件存在性和测试文件属性形成最终候选集合；选择过程只判断文件是否进入集合，不建立候选优先级。规范化结果同时保留判断理由、无效路径、空输出、错误信息和运行状态。

对于 PR $p$，设最终候选集合为 $C_p$，参考集合为 $G_p$，命中集合为 $H_p=C_p\cap G_p$。逐 PR 文件级命中率和召回率定义为：

\[
  \mathrm{HitRate}_p=\frac{|H_p|}{|C_p|},\qquad
  \mathrm{Recall}_p=\frac{|H_p|}{|G_p|}.
\]

实验同时报告总体文件数计算的 micro 指标、逐 PR 平均的 macro 命中率、候选文件数、有效任务数、空输出、无效路径、失败、未完成目标和显式重试。人工复核覆盖准确定位、方法取舍和共同漏检三类案例，并从自然语言质量、代码测试依赖、参考文件规模和运行异常四方面解释结果。

\paragraph{理想结果}

理想情况下，所有原始目标都能够被准确解析为 base 仓库中真实存在的测试文件，不产生无效路径、重复候选、空输出或运行错误；最终候选集合 $C_p$ 与参考集合 $G_p$ 完全一致，即 $C_p=G_p$，从而使文件级命中率和召回率均达到 1。对于参考标签可能遗漏的合理候选，方法应提供足够证据支持人工复核。实际实验将以这一理想结果为参照，分析不同方法在候选范围、参考文件覆盖、运行稳定性和错误类型上的差距。

\section{研究方案}

\subsection{研究方案概述}

本课题以开发任务描述和变更前代码仓库为起点，依次完成测试目标检测、文件级目标选择和效果评价，并根据误检与漏检分析进一步改进定位方法。整体研究流程如图~\ref{fig:research-workflow} 所示。

\begin{figure}[htbp]
  \centering
  \includegraphics[width=\textwidth]{research_workflow.pdf}
  \caption{基于开发任务的测试用例更新目标定位与选择流程}
  \label{fig:research-workflow}
\end{figure}

图中上方主线描述目标定位与选择的完整过程。首先，将 PR 标题、正文和 base 版本仓库构造成统一实验输入；随后，在相同模型、预算、权限和运行环境下，使用三种 SDD 流程及检索对照方法分析开发任务，生成原始测试目标和判断理由；之后，通过路径解析、存在性检查、测试文件判定和去重，将异构输出转换为统一的文件级候选集合；最后，结合命中率、召回率、micro/macro 指标、运行状态和典型案例评价定位效果，并据此分析误检与漏检原因，形成后续方法改进方向。

图中下方支路表示参考答案的隔离过程。head 版本中实际修改的测试文件被独立封存为 gold 集合，不参与任务构造、目标检测和文件选择，只在候选结果固定后进入评价阶段。该设计保证方法无法利用历史修改答案完成定位。评价与案例分析得到的结论再反馈到目标检测阶段，用于改进开发任务表示、代码与测试关联分析以及候选文件选择策略。流程中各项研究内容面临的主要问题及对应方案见表~\ref{tab:research-plan-map}。

\begin{table}[htbp]
  \centering
  \caption{研究内容、主要挑战与拟采用方案}
  \label{tab:research-plan-map}
  \small
  \begin{tabular}{p{2.5cm}p{4.3cm}p{6.2cm}}
    \hline
    研究内容 & 主要问题或挑战 & 拟采用方案 \\
    \hline
    测试更新目标检测 & 自然语言与测试缺少直接对应；参考答案隔离和方法可比性难以保证 & 固定 base 并封存 gold；统一自然语言、模型配置和权限；运行三种 SDD 流程及检索对照 \\
    目标选择与效果评价 & 产物格式和路径口径不一；代理标签不完整且样本规模差异明显 & 统一文件路径解析与选择规则；计算 micro/macro 指标；统计运行状态并复核典型案例 \\
    \hline
  \end{tabular}
\end{table}

\subsection{自然语言驱动的测试更新目标检测方案}
\label{sec:plan-detect-module}

\paragraph{统一检测任务构造}

对每个合格 PR 固定 base SHA，以该提交构建只读仓库和公开任务文件。公开任务只写入仓库标识、PR 编号、base SHA、标题和正文；head 版本代码与测试修改单独生成参考集合并封存，检测容器不挂载 gold 文件。不同方法使用独立工作区，不能读取其他方法的规划产物。任务构造后记录 base 仓库、适配器、镜像和输入摘要，使每次运行能够追溯到确定的版本和配置。

\paragraph{多方法检测流程}

OpenSpec、Spec Kit 和 cc-sdd 接收同一任务，按照各自流程生成规格、设计和任务产物。统一提示要求方法只分析变更需求、检查 base 仓库并列出可能需要更新的已有测试文件及理由，不执行实现、不修改仓库、不访问 head。随后收集规划产物中的原始测试路径、测试符号、判断依据和运行证据。

实验统一使用 DeepSeek-V4.1-Flash，配置标识为 \texttt{deepseek-flash}，推理强度为 high。每次运行的 rollout token 上限为 100,000 tokens，时间上限为 2,700 秒。一个“PR--方法”组合构成一个检测目标，每个阶段计划执行一次；若因上游服务中断等非方法因素失败，则使用显式重试恢复，并单独记录尝试次数和失败原因，不把重试作为独立重复实验。

对照方法拟适配 IssueExec 的测试检索阶段\cite{liu2026issueexec}。该方法使用 PR 标题和正文替代 issue 文本，在 base 测试集中执行 BM25 检索，并在条件允许时由大语言模型筛选相关测试。若加入历史知识增强，只使用 base 提交以前的记录，避免引入未来信息。

\paragraph{检测结果与过程产物}

检测过程为每个“PR--方法”目标生成一组原始测试目标及其判断理由，并保存规格、设计、任务、运行日志和状态记录。这些产物既是后续文件级选择的输入，也是分析不同方法如何理解同一自然语言需求的证据。

\subsection{文件级目标选择与效果评价方案}
\label{sec:plan-select-module}

\paragraph{候选文件解析与选择}

建立统一结果解析器处理三种 SDD 产物和检索结果。解析器先识别文件路径、测试类或测试函数，再将绝对路径、相对路径和符号引用转换为仓库相对文件路径；随后检查目标是否存在于 base 仓库、是否属于测试文件，并删除重复项。无法解析或不存在的路径记录为无效路径，不计入有效候选。规范化结果包含任务标识、方法、候选文件集合、判断理由、运行状态和错误信息。

\paragraph{指标计算与状态统计}

在所有检测结束后解封参考集合，按文件路径计算 $C_p$ 与 $G_p$ 的交集、命中率和召回率。micro 指标通过汇总全部文件数计算，macro 命中率通过逐 PR 计算后取平均。运行层面分别统计有效输出、空输出、无效路径、上游中断、错误和显式重试，避免把运行条件问题与方法判断结果混合解释。

\paragraph{案例复核与结论形成}

选择候选与 gold 完全一致的准确案例、不同方法表现出候选范围取舍的案例，以及三种方法共同漏检的案例进行人工复核。对于多文件 PR，检查遗漏是否集中在共享基类、间接依赖或跨模块测试；对于未出现在 gold 中但具有合理关联的候选，结合代码和测试证据说明代理标签的边界。最终从候选准确性、参考文件覆盖、运行稳定性和可解释性四方面比较方法。

\subsection{实验控制与有效性保障}

实验的自变量为检测方法，因变量包括候选数、命中数、命中率、召回率和运行状态。目标检测过程统一 PR 样本、base 提交、自然语言、模型、推理强度、token 预算、超时、镜像、任务提示和工具权限；文件选择与评价过程统一路径解析、文件检查、去重规则和指标实现。正式实验将先恢复 4 个上游中断目标和 PR 7985 的 cc-sdd 错误，再生成统一汇总。

研究主要面临三类有效性威胁。历史测试修改可能遗漏理论上也应更新的文件，因此需要人工检查典型误检；模型服务中断可能被错误解释为方法漏检，因此运行异常单独报告；当前结果来自 Scrapy，不能直接推广到所有 Python 项目，因此后续扩展项目时需要重新检查自然语言分布和阈值适用性。

\section{进度安排，预期达到的目标}

\subsection{已完成的工作}

\paragraph{PR 样本筛选与质量检查}

项目已从 ansible、matplotlib、pandas 和 scrapy 四个 Python 项目采集 PR 元数据、版本标识、CI 状态和文件变更。第一轮通过 \texttt{merge\_quality} 和 \texttt{code\_presence} 检查合并质量与代码存在性，得到 938 个预筛选 PR；第二轮检查 CI，排除明确失败或状态不可确认的样本；随后检查自然语言质量、测试修改、base/head 完整性、仓库可复现性和重复样本。

自然语言质量使用 PR 标题、正文和提交消息的有效词数。根据词数的对数变换分布、分位数和最低信息量，标题、正文和提交消息阈值分别取 4、14 和 6，联合规则为：

\[
  \mathrm{InfoQuality}(p)=
  \bigl(T_p\geq4\bigr)\land
  \bigl((B_p\geq14)\lor(M_p\geq6)\bigr).
\]

其中 $T_p$、$B_p$ 和 $M_p$ 分别为标题、正文和提交消息有效词数。提交消息只作为筛选证据，不加入当前 detect 输入。筛选过程保留各字段词数、判定结果和剔除原因。目前已完成 30 个 Scrapy PR 的筛选与质量检查，这些 PR 按 PR 累计包含 145 个参考测试文件。

\paragraph{检测任务与初步实验}

已为 30 个 PR 构造隔离的 base 检测任务，并分别运行 OpenSpec、Spec Kit 和 cc-sdd，共形成 90 个检测目标。汇总记录中 89 行执行状态为 OK、1 行为 ERROR。其中 Spec Kit 有 4 个目标因上游 token 池中断而未产生候选。为保持 90 个目标的总体统计完整性，当前描述性结果将这 4 个目标按未命中计入；同时将其标记为运行异常，正式方法比较前需要恢复运行并重新计算指标。表~\ref{tab:pilot-detect-results} 给出包含上述异常目标的当前整体文件级结果，micro 召回率以 145 个参考测试文件为共同分母。

\begin{table}[htbp]
  \centering
  \caption{30 个 Scrapy PR 的初步文件级检测结果}
  \label{tab:pilot-detect-results}
  \small
  \begin{tabular}{lcccc}
    \hline
    方法 & 命中/候选 & micro 命中率 & macro 命中率 & micro 召回率 \\
    \hline
    OpenSpec & 56/136 & 0.412 & 0.623 & 0.386 \\
    Spec Kit & 46/90 & 0.511 & 0.582 & 0.317 \\
    cc-sdd & 58/166 & 0.349 & 0.588 & 0.400 \\
    \hline
  \end{tabular}
\end{table}

Spec Kit 的候选文件最少，micro 命中率最高；OpenSpec 的逐 PR 平均命中率最高；cc-sdd 找回的参考文件数最多，同时候选文件也最多。这些差异说明三种方法在候选范围与参考文件覆盖方面存在不同取舍，仍需结合具体案例解释。需要说明的是，表中 Spec Kit 的结果包含 4 个上游 token 池中断目标；这些目标在当前描述性统计中按未命中计入，会压低其召回表现。该表用于呈现目前全部 90 个目标的运行结果，不能作为排除运行异常后的正式方法比较结论。

\paragraph{典型案例人工复核}

\begin{enumerate}
  \item \textbf{PR 7918：三种方法均准确定位。}该 PR 将 \texttt{ScrapyJSONEncoder} 的日期和时间序列化方式改为 ISO 8601，参考文件为 \texttt{tests/test\_exporters.py} 和 \texttt{tests/test\_utils\_serialize.py}。三种方法均只输出这两个文件，结果为 2/2。自然语言、生产代码和测试断言之间具有明确对应关系。
  \item \textbf{PR 7846：召回能力与候选范围存在取舍。}该 PR 为 Twisted 的 IDNA 辅助函数增加兼容处理。OpenSpec、Spec Kit 和 cc-sdd 的结果分别为 1/5、1/7 和 2/8。cc-sdd 找回两个参考文件，同时引入更多语义相关候选，说明扩大依赖分析范围能够降低漏检，也会增加人工核查成本。
  \item \textbf{PR 7966：三种方法共同漏检。}该 PR 标题为“Address warnings from the VCS workflow”，正文为空，参考文件是 \texttt{tests/test\_utils\_console.py}。三种方法分别给出 14、9 和 11 个候选文件但均未命中，说明缺少具体警告类型和受影响组件时，方法容易根据仓库常见关联产生系统性偏离。
\end{enumerate}

PR 7808、7841 和 7968 的参考文件数分别为 41、28 和 28，合计占全部参考文件的 97/145，多文件 PR 是当前漏检的重要来源。后续还需处理 4 个上游中断目标和 PR 7985 的 cc-sdd 运行错误，再完成正式比较。

\subsection{进度安排}

根据目前工作进展，研究工作按开题、中期和结题三个阶段安排如下：

\begin{itemize}
  \item \textbf{开题检查前（2026 年 9 月至 10 月）}：查阅测试维护、代码测试共同演化、SDD 及 IssueExec 等文献，明确自然语言驱动的文件级检测问题；已完成 30 个 Scrapy PR 的样本筛选与质量检查，取得三种 SDD 流程的初步检测结果，明确文件级评价指标、运行异常和实验方案，完成开题报告。
  \item \textbf{开题至中期前（2026 年 10 月至 2027 年 3 月）}：在已筛选样本上完善参考答案隔离与统一评价流程；统一实验条件并恢复受环境差异或运行异常影响的目标；实现或适配 IssueExec 的测试检索对照方法，形成可比较的中期结果。
  \item \textbf{中期至结题前（2027 年 3 月至 5 月）}：完成正式检测实验、文件级指标统计和典型误差复核，分析方法的适用范围与局限；撰写、修改毕业论文，整理实验材料并完成结题准备。
\end{itemize}

\subsection{预期目标}

形成字段边界明确、版本可追溯的 PR 测试文件检测样本；建立从自然语言到 SDD 规划产物、再到候选测试文件的可复现实验流程；获得三种 SDD 流程及测试检索对照方法在统一条件下的文件级结果；总结自然语言质量、代码测试关联和修改规模对定位效果的影响，为后续测试更新研究提供数据与方法基础。

\section{课题已具备和所需的条件、经费}

现有仓库已具备 PR 采集、base/head 版本记录、自然语言质量分析、Python/pytest 测试提取、任务隔离、SDD 基准运行和文件级指标汇总代码，并已获得 30 个 Scrapy PR 的阶段性结果。正式实验仍需统一的镜像环境、稳定的模型运行条件、足够的 API 调用资源，以及人工复核标签的时间。课题预计经费为 5,000 元，全部用于购买 DeepSeek-Flash 模型实验所需的 API token（调用额度）。

\section{研究过程中可能遇到的困难和问题，解决的措施}

\begin{enumerate}
  \item \textbf{参考标签不完整。}历史 PR 修改的测试文件不一定覆盖全部应更新目标。拟抽取典型误检进行人工复核，将相关但未修改的文件单独讨论，并明确代理标签边界。
  \item \textbf{自然语言质量不稳定。}部分 PR 标题笼统、正文为空，难以提供具体组件和行为线索。拟保留词数筛选结果，按信息完整程度分析错误，并研究提交消息作为辅助输入的增量价值。
  \item \textbf{多文件 PR 影响总体指标。}少数大规模修改包含大量参考文件，可能主导 micro 结果。拟同时报告 micro 和 macro 指标，并按参考文件规模进行分组分析。
  \item \textbf{比较和复现困难。}部分目标存在上游服务中断、显式重试和运行错误。拟固定镜像、模型、预算、超时和路径解析规则，恢复受影响任务，并完整报告成功与失败状态。
\end{enumerate}

\section{主要参考文献}

% 当前 ../reference.bib 含演示条目，开题报告使用单独核实的文献库。
\nocite{fatima2024flakyfix,guo2022unixcoder,jang2025unitcon,liu2025cote,alshahwan2024testgenllm,schaefer2024empirical}
\bibliographystyle{hithesis}
\bibliography{reference_opening}

% Local Variables:
% TeX-master: "../report"
% TeX-engine: xetex
% End:
